\documentclass[11pt]{article}
\usepackage[margin=25mm]{geometry}
\usepackage{fontspec}
\setmainfont{Latin Modern Roman}
\usepackage{amsmath,amssymb,mathtools}
\usepackage{unicode-math}
\setmathfont{Latin Modern Math}
\usepackage{xurl,hyperref}
\hypersetup{colorlinks=true,urlcolor=blue}
\setcounter{secnumdepth}{0}
\setlength{\parindent}{0pt}
\setlength{\parskip}{5pt}
\setlength{\emergencystretch}{3em}
\newcommand{\sref}[2]{\hyperlink{source:#1:#2}{[S#2]}}
\newcommand{\cataloguescope}[1]{\paragraph{Recorded target.}#1}
\begin{document}
\section{Maximum Givens in Minimal Sudoku}\label{statement}

\subsection{Definitions and mathematical statement}

\paragraph{Shared notation.}$\mathbb N=\{1,2,\ldots\}$, and $\mathbb Z,\mathbb Q,\mathbb R,\mathbb C$ denote the integers, rationals, reals and complex numbers. Any other conventions are specified locally.


Let $C=\{1,\ldots,9\}^2$ be the cells of a standard $9\times9$ Sudoku. A solution is a map $s:C\to\{1,\ldots,9\}$ whose values in each row, each column and each of the nine $3\times3$ boxes are exactly $1,\ldots,9$. A puzzle is a partial map $p:D\to\{1,\ldots,9\}$ with $D\subseteq C$; its givens are the $|D|$ assigned cells. Write $\operatorname{Comp}(p)=\{s:\ s\text{ is a solution and }s|_D=p\}$. A puzzle is minimal if $|\operatorname{Comp}(p)|=1$ and, for every $d\in D$, deleting that clue gives $|\operatorname{Comp}(p|_{D\setminus\{d\}})|\ge2$.
\paragraph{Statement recorded in the supplied source.}
Determine
\[
 M=\max\bigl\{|D|:\ p:D\to\{1,\ldots,9\}\text{ is a minimal standard Sudoku puzzle}\bigr\}.
\]
Minimality means that every given is essential to uniqueness; it does not mean having the fewest givens among all puzzles. \sref{statement}{1}\sref{statement}{2}
\subsection{Short English statement}
What is the largest number of givens that a standard 9 by 9 Sudoku can have while every individual given is essential to its unique solution?
\subsection{Sources}
\begin{description}
\item[\hypertarget{source:statement:1}{S1}] ULAM.AI and the UnsolvedMath Contributors, \emph{UnsolvedMath: A Curated Collection of Open Mathematics Problems}, catalogue entry. CC BY 4.0. \url{https://huggingface.co/datasets/ulamai/UnsolvedMath}.
\item[\hypertarget{source:statement:2}{S2}] Gary McGuire, Bastian Tugemann and Gilles Civario, \emph{There is no 16-Clue Sudoku: Solving the Sudoku Minimum Number of Clues Problem via Hitting Set Enumeration} (2013), arXiv:1201.0749v2, Section 1 (Sudoku rules and unique completion), and Section 5.1 (clue sets and unavoidable sets). \url{https://arxiv.org/abs/1201.0749v2}.
\end{description}
\label{end:statement}
\end{document}
